\section{Goal Classifiers: aprender la recompensa a partir de ejemplos de objetivos}

\subsection{Idea básica}

Se recolectan estados exitosos (ejemplos positivos) y estados fallidos (ejemplos negativos), se entrena un clasificador binario y se usa su salida como señal de recompensa.

\begin{enumerate}
    \item Recolectar el conjunto de estados exitosos $D^+$ y el conjunto de estados fallidos $D^-$
    \item Entrenar un clasificador binario: entrada $s_i$, etiqueta $\mathbf{1}(s_i \in G)$
    \item Usar la salida del clasificador como recompensa para RL
\end{enumerate}

\begin{warningbox}{Reward Hacking}
El algoritmo de RL buscará activamente estados que hagan que el clasificador produzca una puntuación alta. Puede encontrar estados con los que el clasificador \textbf{nunca se entrenó}, es decir, explotar las debilidades del clasificador en lugar de completar realmente la tarea.
\end{warningbox}

\subsection{Entrenamiento adversario para resolver reward hacking}

Idea de solución: agregar los estados visitados por RL al conjunto de ejemplos negativos, obligando al clasificador a actualizarse continuamente para evitar ser explotado.

\begin{enumerate}
    \item Inicializar $D^+$ (estados exitosos) y $D^-$ (estados fallidos)
    \item Entrenar el clasificador con $D^+$ y $D^-$ (conjunto de datos balanceado)
    \item Recolectar experiencia $s_t, a_t, \ldots$ con la política $\pi$
    \item Actualizar la política con la recompensa del clasificador
    \item Agregar los estados visitados a los ejemplos negativos: $D^- \leftarrow D^- \cup \{s_t\}$
    \item Repetir
\end{enumerate}

\begin{knowledgebox}{Por qué funciona este método}
El clasificador ya no puede ser explotado. Pregunta clave: ¿qué pasa si la política realmente tiene éxito? Mientras se mantenga el muestreo balanceado de $D^+$ y $D^-$, el clasificador seguirá produciendo $p \geq 0.5$ para los estados verdaderamente exitosos.
\end{knowledgebox}

Los experimentos muestran que, en tareas de robótica, un goal classifier entrenado con 50 demostraciones combinado con RL alcanzó una tasa de éxito del 62\%, muy superior al 26\% del aprendizaje por imitación directo.

\subsection{Resumen de la sección}

Goal classifier es un método simple y eficaz de aprendizaje de recompensa. El entrenamiento adversario es clave para evitar el reward hacking.
